\documentclass[11pt]{article}
\usepackage{xcolor}
\definecolor{darkgreen}{rgb}{0.0,0.6,0.0}

\usepackage{amsmath,amssymb,amsthm}
\usepackage{proof}
\usepackage{verbatim}
\usepackage[backend=bibtex,style=alphabetic]{biblatex}
\usepackage{fancyhdr}
\usepackage{titlesec}
\usepackage{microtype}
\usepackage[margin=1in]{geometry}

\pagestyle{fancy}
\fancyhf{}
\rhead{Chuta Sano}
\lhead{Teaching Statement}
\rfoot{\thepage}
\setlength{\headheight}{13.59999pt}

\titleformat{\section}{\large\bfseries}{}{0em}{}[\titlerule]

\begin{document}

\title{Teaching Statement}
\author{Chuta Sano}
\maketitle{}

In academia, teaching, advising, and research work hand in hand, complementing one another.
I value teaching and advising not only for their own sakes but also for the many unexpected learning benefits from engaging with students of diverse academic backgrounds.

At McGill University, I served as a teaching assistant for \emph{Programming Languages and Paradigms}, \textbf{an undergraduate-level functional programming course with about 300 students} and \emph{Logic and Computation}, \textbf{a graduate-level theory course with about 40 students}.
I was awarded the Best Teaching Assistant Award in 2023 based on student nominations for my work in \emph{Logic and Computation}.
I am currently \textbf{advising two undergraduate students} on their honors-level theses and have informally advised several other student groups on class projects as part of \emph{Logic and Computation}.
These experiences have helped me develop a flexible teaching and advising style that adapts to the needs of different students and course logistics.

In the following, I describe my general teaching philosophy and elaborate on my teaching experiences. I then describe my advising philosophy and experiences.

\section*{Teaching}
Computer science is a broad field that encompasses both theory and software-engineering practices, each requiring different teaching styles.
Furthermore, logistical details such as class size, classroom setup, and available resources also influence what teaching styles are most effective.
I emphasize adapting my teaching style primarily based on two factors: whether the course topic is in theory or software engineering, and whether the class size is small ($\sim 50$ max) or large.
In the interest of brevity, I will consider two categories, \emph{a theory course in a small class} and \emph{a programming course in a large class}.
These two, in my experience, cover the majority of scenarios and are therefore the most important to discuss.

\subsection*{Theory Course in a Small Class}
For any given lecture, my general workflow consists of \textbf{preparing lecture notes, delivering the lecture on a blackboard, and releasing my notes} after adjusting them based on the actual lecture.
Generally, I find this style of teaching more effective than using pre-prepared slides since it allows for more flexibility and interactivity during the lecture.

When introducing new theoretical concepts, I focus on \textbf{starting small to build intuition} before giving formal definitions.
For example, in \emph{Logic and Computation}, I gave several lectures covering Linear Logic, which requires that assumptions are used exactly once.
I started by giving simple examples of resource usage in everyday life, such as the concept of purchasing items with money.
I then gradually formalized this idea to build a logical system that models shopping where assumptions correspond to money and purchased items correspond to conclusions.
This way, I was able to motivate the restrictions of Linear Logic based on the money analogy: the absence of contraction represents the inability to duplicate money, and the absence of weakening represents the inability (or undesirability) to throw away money.
This made the transition to Linear Logic much smoother, as students gained intuition on the idea of treating assumptions as resources and why this restriction is useful.

My lectures were very well-received by students; students acknowledged it during their nomination of me for the Best Teaching Assistant Award, where they noted, ``The lectures you gave during the semester were fantastic''.
They also highlighted, ``Your patience, kindness, and clear explanations of difficult topics'', which I believe is a testament to my teaching style.

I also believe a core part of student learning, especially for theory courses, is through reflecting on their assignments.
In Logic and Computation, I spent a significant amount of effort in \textbf{grading assignments to provide clear and personalized feedback}, and our policy to allow resubmissions of assignments gave students a strong incentive to take my feedback seriously.
The positive effect of this was clear; students attended office hours much more frequently than in past semesters, often bringing in their previous submissions and my feedback to discuss how they could improve.
This was also reflected in student feedback, where they noted
``The extra effort you made to explain complex topics going beyond the material in class''.


\begin{comment}
I also place a heavy emphasis on designing \textbf{assignments that encourage students to explore the material on their own}.
For example, in \emph{Logic and Computation}, I designed an assignment on the aforementioned Linear Logic, where I asked students to explore the role of the restrictions of Linear Logic by making explicit the structural rules of weakening and contraction.
In particular, I asked students to prove propositions both with and without these structural rules to see for themselves how the presence or absence of these rules affected the provability of certain statements.
\end{comment}


\subsection*{Programming Course in a Large Class}
The aforementioned style of giving flexible lectures via blackboard and personalized feedback on assignments does not scale well to large class sizes.
Thus, I value giving organized and well-structured lectures through a computer-based setup, both through slides and live-coding demonstrations.
I still retain some flexibility and interactivity with the slide-based presentation through the use of a drawing tablet attached to my computer, allowing me to seemlessly write on slides.
Broadly speaking, I teach particular programming concepts and techniques by introducing small examples that serve as motivating cases.
I then gradually build on these examples to arrive at the core concept and then give additional applications of the concept to solidify student understanding.

However, large courses especially require iterating on the course materials and structure over multiple semesters, due to the complexity arising from the large number of students and TAs/staff involved.
Therefore, I will spend a significant amount of time understanding and adapting to existing course structures and materials from past semesters.
For example, while I believe the web-based LearnOCaml platform is a great tool to teach functional programming (and OCaml), I would first familiarize myself with any language and tools that past courses used and consider trade-offs of different design choices in the course before making any changes.

As a teaching assistant for \emph{Programming Languages and Paradigms} (2022 Fall), I was responsible for designing and grading assignments, grading exams, and holding office hours.
Our class used OCaml as the main programming language and a web-based platform LearnOCaml for assignments, which included an interpreter and an auto-grader.
Thus, the development of assignments involved designing code templates and writing (randomized) test cases for the auto-grader in addition to the usual assignment specifications.
Given the large class size, I placed a heavy emphasis on \textbf{designing clear and unambiguous assignment specifications} -- I found that this helps not only to reduce the number of confusion amongst students but also to increase the number of students who can help others.
While I actively monitored the course forum (via Discord), I noticed the assignments I designed had much more interaction between students.

\section*{Advising}
I view advising as an opportunity to mentor students in their academic and professional growth.
There is, unfortunately, no one-size-fits-all approach to advising, as each student has different strengths and work-styles;
some students prefer to dictate the pace of their project, while others prefer more structured guidance with frequent check-ins.
These kinds of expectations are tuned over time, so I do my part by assuring students that I am available at any time, and that they can approach me with any questions or concerns.


\subsubsection*{Experience}
I have advised several student groups on class projects as part of \emph{Logic and Computation}, where I provided feedback on their project progress, presentation, and final report.
I am also currently advising two undergraduate students on their honors-level theses; both students are working on projects based on my research, and I give regular feedback on their progress and help them navigate technical challenges.


\section*{Teaching Plan}
I believe I am capable of teaching just about any core computer science course.
However, given my background, I am especially well-suited to teach courses in general programming with rigor (such as CSC110Y1 or CSC111H1) compilers (CSC488H1), and theory of computation (CSC236H1).
I would also like to teach graduate-level courses on topics that are aligned with my research, such as verification and foundations of programming languages.

\begin{comment} Generic
I believe I am capable of teaching just about any core computer science course.
However, given my background, I am especially well-suited to teach courses in general programming (especially functional programming), compilers (COMS W4115), and theory of computation.
I would also like to teach graduate-level courses on topics that are aligned with my research, such as verification and foundations of programming languages.
\end{comment}

\begin{comment} Cornell Ithaca
However, given my background, I am especially well-suited to teach courses in general programming (especially CS 3110), compilers (CS 4120), and theory of computation (CS 4110).
I would also like to teach graduate-level courses on topics that are aligned with my research, such as verification and foundations of programming languages.
\end{comment}

\begin{comment} Cornell Tech
While I believe my background demonstrates that I am capable of teaching just about any core computer science course, I am also well-suited to teach graduate courses on topics that are aligned with my research, such as verification and foundations of programming languages.
I am also interested in designing and teaching new courses based on ongoing research, such as a course on formal verification techniques for programming languages.
\end{comment}

\begin{comment} UCB
However, given my background, I am especially well-suited to teach courses in general programming (especially CS 61A) and compilers (CS 164).
I would also like to teach graduate-level courses on topics that are aligned with my research, such as verification and foundations of programming languages.
\end{comment}

\begin{comment} U of T
I believe I am capable of teaching just about any core computer science course.
However, given my background, I am especially well-suited to teach courses in general programming with rigor (such as CSC110Y1 or CSC111H1) compilers (CSC488H1), and theory of computation (CSC236H1).
I would also like to teach graduate-level courses on topics that are aligned with my research, such as verification and foundations of programming languages.
\end{comment}

\end{document}

